\chapter{分布与弱导数}

\begin{solution}{distribution-regular-continuity}
若 \(\operatorname{supp}\varphi\subset K\Subset\Omega\)，则
\[
 |\langle T_f,\varphi\rangle|
 \le\int_K|f||\varphi|
 \le\|f\|_{L^1(K)}\|\varphi\|_\infty.
\]
所以局部估计成立，阶数可取零。若 \(\varphi_n\to0\) 于试验函数空间，
存在共同 \(K\) 且上确界范数趋零，上式给 \(T_f\varphi_n\to0\)。
\end{solution}

\begin{solution}{distribution-dirac-not-regular}
反设 \(\delta_0=T_f\)。取 \(\eta\in C_c^\infty(B(0,1))\) 且
\(0\le\eta\le1,\eta(0)=1\)，令 \(\eta_n(x)=\eta(nx)\)。则
\(\delta_0(\eta_n)=1\)，而
\[
 \left|\int f\eta_n\right|
 \le\int_{B(0,1/n)}|f|\to0
\]
由局部积分的绝对连续性成立，矛盾。一般 \(x_0\) 只需平移。
\end{solution}

\begin{solution}{distribution-derivative-continuity}
若 \(\operatorname{supp}\varphi\subset K\)，对 \(T\) 使用阶数 \(m\) 的估计，
\[
 |\langle\partial^\alpha T,\varphi\rangle|
 =|\langle T,\partial^\alpha\varphi\rangle|
 \le C_K\max_{|\beta|\le m}
 \|\partial^{\alpha+\beta}\varphi\|_\infty.
\]
这是阶数至多 \(m+|\alpha|\) 的局部估计，故导数连续。
\end{solution}

\begin{solution}{distribution-heaviside}
正文已得 \(H'=\delta_0\)。函数 \(|x|\) 局部绝对连续，故其弱导数为
\(\operatorname{sgn}x\)。而
\(\operatorname{sgn}x=2H-1\)，所以
\[
 D|x|=\operatorname{sgn}x,\qquad D^2|x|=2\delta_0.
\]
\end{solution}

\begin{solution}{distribution-product-rule}
对试验函数 \(\varphi\)，
\[
 \langle\partial_j(\psi T),\varphi\rangle
 =-\langle T,\psi\,\partial_j\varphi\rangle
 =-\langle T,\partial_j(\psi\varphi)\rangle
   +\langle T,(\partial_j\psi)\varphi\rangle.
\]
右侧等于
\(\langle\psi\,\partial_jT+(\partial_j\psi)T,\varphi\rangle\)。
\end{solution}

\begin{solution}{distribution-piecewise-jumps}
先截去分点附近的小区间，在各剩余段上分部积分，再令截去的长度趋零。
题设中的 \(f'_{\mathrm{cl}}\in L^1_{\mathrm{loc}}\) 保证内部导数积分收敛；
分点两侧的边界项由单侧极限合并为
\((f(a_j+)-f(a_j-))\varphi(a_j)\)，即相应 Dirac 项。
在实线上定义周期锯齿函数
\(f(x)=x-2\pi k\)，其中 \((2k-1)\pi<x<(2k+1)\pi\)。
它在内部导数为一，在每个拼接点从 \(\pi\) 跳到 \(-\pi\)，故
\[
 Df=1-2\pi\sum_{k\in\mathbb Z}\delta_{(2k+1)\pi}.
\]
和式局部有限：每个紧支撑试验函数只遇到有限个分点。因此它在
\(\mathcal D'(\R)\) 中有定义，无须另行引入圆周上的分布。
\end{solution}

\begin{solution}{distribution-weak-unique}
若 \(g_1,g_2\) 都是弱导数，则
\(\int(g_1-g_2)\varphi=0\) 对所有试验函数成立。正则分布嵌入的单射性给
\(g_1=g_2\) 几乎处处。
\end{solution}

\begin{solution}{distribution-weak-ac-one-d}
令 \(G(x)=\int_a^xg\)，则 \(G\) 局部绝对连续且弱导数为 \(g\)。
\(h=f-G\) 的弱导数为零。取正长度紧区间 \(J\Subset I\)，选
\(J\subset U\Subset I\)，其中 \(U\) 为开区间；令 \(\widetilde h\)
为 \(h\mathbf1_U\) 的零延拓。对充分小的 \(\varepsilon\) 及 \(x\in J\)，
\[
 (\widetilde h*\rho_\varepsilon)'(x)
 =-\int_U h(y)\partial_y\rho_\varepsilon(x-y)\,dy=0,
\]
因为该核在 \(y\) 变量上是允许的试验函数。这直接从零弱导数的定义推出，
无需先引用后文的分布卷积交换公式。故卷积在 \(J\) 上为常数 \(c_\varepsilon\)，并在 \(L^1(J)\) 中趋于
\(h\)，常数列因其 \(L^1\) 距离趋于零而为 Cauchy，故 \(h=c\) 几乎处处。
用递增正长度紧区间穷尽 \(I\)，重叠区间上的常数相同，可数并例外集仍零测，
得到全局局部绝对连续代表 \(G+c\)。反向由绝对连续
分部积分直接成立。
\end{solution}

\begin{solution}{distribution-mollifier-commute}
对 \(x\) 求导时，
\[
 \partial_x^\alpha\langle T,\rho_\varepsilon(x-\cdot)\rangle
 =\langle T,\partial^\alpha\rho_\varepsilon(x-\cdot)\rangle.
\]
按分布导数定义把导数移到 \(T\) 上，得到交换公式。对试验函数 \(\varphi\)，
取 \(\varepsilon\) 足够小，使试验函数支撑在 \(\Omega_\varepsilon\) 中。
按正文证明，用具有共同紧支撑、各阶导数一致收敛的 Riemann 和交换配对，
得到
\[
 \langle T*\rho_\varepsilon,\varphi\rangle
 =\langle T,\widetilde\rho_\varepsilon*\varphi\rangle.
\]
右侧试验函数有共同紧支撑且各阶导数一致趋于 \(\partial^\alpha\varphi\)，
分布连续性给收敛。这里左侧是 \(\Omega_\varepsilon\) 上的积分配对，
并未将正则化定义成整个 \(\Omega\) 上的函数；每个固定试验函数最终可配对即足够。
\end{solution}

\begin{solution}{distribution-fundamental-lemma}
在任意球 \(B\Subset\Omega\) 中延拓 \(f\mathbf1_B\) 为零，并与磨光核卷积。
题设给
\[
 (f*\rho_\varepsilon)(x)
 =\int f(y)\rho_\varepsilon(x-y)dy=0
\]
对离边界足够远的 \(x\) 成立。另一方面磨光近似在 \(L^1_{\rm loc}\) 中趋于
\(f\)，故 \(f=0\) 于每个较小球中几乎处处。用可数个这样的球覆盖
\(\Omega\)，得到结论。
\end{solution}

\begin{solution}{distribution-x-delta}
\[
 \langle x\delta_0,\varphi\rangle=0.
\]
又
\[
 \langle x\delta_0',\varphi\rangle
 =\langle\delta_0',x\varphi\rangle
 =-(x\varphi)'|_{0}=-\varphi(0),
\]
故 \(x\delta_0'=-\delta_0\)。最后
\[
 \langle x^2\delta_0'',\varphi\rangle
 =(x^2\varphi)''|_0=2\varphi(0),
\]
所以 \(x^2\delta_0''=2\delta_0\)。
\end{solution}

\begin{solution}{distribution-fundamental-solution-one-d}
由 \(D|x|=\operatorname{sgn}x\) 和
\(D\,\operatorname{sgn}x=2\delta_0\)，得
\((|x|/2)''=\delta_0\)。若 \(f\in C_c^\infty\)，令
\[
 u=E*f,\qquad E(x)=|x|/2.
\]
本题无需先定义任意两个分布的卷积。由于 \(f\) 紧支撑，下列积分绝对收敛：
\[
 u(x)=\frac12\int_{-\infty}^x(x-y)f(y)\,dy
      +\frac12\int_x^\infty(y-x)f(y)\,dy.
\]
积分上限求导时，两项活动端点的因子均为零，所以
\[
 u'(x)=\frac12\int_{-\infty}^x f(y)\,dy
       -\frac12\int_x^\infty f(y)\,dy,\qquad u''(x)=f(x).
\]
因此 \(u\) 是经典解，实际上为光滑函数。任意两个经典解之差二阶导数为零，
故全部经典解为 \(E*f+ax+b\)，其中 \(a,b\in\mathbb K\)。
\end{solution}

\begin{solution}{distribution-no-general-product}
取非负标准核 \(\rho\in C_c^\infty(-1,1)\)，\(\int\rho=1\)，令
\(h_\varepsilon=H*\rho_\varepsilon\)、
\(a_\varepsilon(x)=h_\varepsilon(x-2\varepsilon)\)、
\(b_\varepsilon(x)=h_\varepsilon(x+2\varepsilon)\)。三者光滑且介于0与1之间，
在每个非零点都趋于 \(H\)，故由紧集上的控制收敛，
\(T_{a_\varepsilon},T_{b_\varepsilon}\to T_H\)。另一方面，换元给
\[
 \int\rho_\varepsilon(x)\varphi(x)\,dx
 =\int\rho(z)\varphi(\varepsilon z)\,dz\longrightarrow\varphi(0),
\]
故 \(T_{\rho_\varepsilon}\to\delta_0\)。但在
\(\operatorname{supp}\rho_\varepsilon\subset[-\varepsilon,\varepsilon]\) 上，
\(a_\varepsilon=0\)、\(b_\varepsilon=1\)，于是
\[
 a_\varepsilon\rho_\varepsilon=0,\qquad
 b_\varepsilon\rho_\varepsilon=\rho_\varepsilon.
\]
令 \(\varepsilon=1/n\)，若题设的 \(P\) 存在，经典乘法的一致性与联合序列
连续性使左一列的极限为 \(P(T_H,\delta_0)=0\)，左二列又给
\(P(T_H,\delta_0)=\delta_0\)，矛盾。结论针对所列两项要求；
它不否定特定分布之间另有定义的乘积，也不声称任意无连续性要求的代数规则都不存在。
\end{solution}
